% sections/discussion.tex
\section{Discussion}\label{sec:discussion}

\subsection{What the simulation adds to the synthesis}

The review found that association models fitted to contemporaneous, linearly specified exposure dominate the literature. The simulation establishes what that habit costs under conditions that resemble the reviewed data: annual site by year panels, counts with overdispersion, and a true response that is delayed and mildly saturating. The naive GLM did not merely lose precision; it returned estimates close to zero when the true lag was one year or more, which in an applied study would read as no detectable effect. A lag aware GLM recovered the effect when the lag it assumed was correct and failed when it was not, which is the practical problem, because the lag is seldom known in advance. The distributed lag model, which estimates the lag profile rather than assuming it, was the only class whose point estimates did not collapse under any lag condition; it paid for that robustness in precision, and its intervals, while covering the truth an order of magnitude more often than the linear models', still fell short of nominal. This ordering does not depend on region, grid size or spatial range, which is the simulation's answer to the question of whether its findings transfer: the mechanism is the mismatch between the model's temporal assumption and the biology's, and that mismatch is indifferent to where the temperature series came from.

The interval results deserve emphasis because they are the least intuitive. Linear models returned standard errors that were correctly computed for the model fitted; their intervals failed because the point estimate was biased by curvature the model could not represent, not because the uncertainty was understated. A reader who sees a narrow, well behaved confidence interval from a GLM has no way to tell, from the interval alone, that it is centred on the wrong value. Only a model that can represent the curvature, or a diagnostic that detects it, reveals the problem. This is why item 6 of the reporting standards is rated critical, and why the workflow asks for a smooth term to be tested even when a linear term is ultimately reported.

The inference versus prediction result is a reminder that model selection criteria encode a purpose. Averaged across the design, the model that predicted held out years best was the naive GLM, whose coefficient was the least trustworthy at non zero lags, and the random forest, which returns no coefficient at all, predicted competitively throughout. Fisheries applications typically need both properties, for attribution and for management advice respectively, and the simulation shows that no single model delivers both. Reporting one model of each class is a cheap safeguard.

\subsection{Scope and limitations of the simulation}

The simulation is a statement about model behaviour under its design, and the design was chosen to be transparent rather than comprehensive. Three limitations bound its reach. The data generating process is a single functional form with one source of non linearity and no interaction between exposure and site; real responses may saturate differently, respond to peak rather than cumulative intensity, or vary with life stage. The panel structure fixes an annual response cadence; studies with seasonal or monthly responses face a different and generally easier lag estimation problem. And the spatial structure enters only through correlated noise in the exposure field, not through spatially varying effects, so the simulation can show whether a spatial random field model recovers a fixed effect reliably but cannot show the conditions under which spatial variation in the effect itself becomes estimable. Each of these is a design choice recorded in a single configuration file, and the repository exists so that readers who find the choices unrepresentative of their system can change them.

The sdmTMB results carry an additional qualification. The model was fitted on \sdmtmbPanelsN{} of the \nPanels{} exposure panels (\sdmtmbScenarioFrac{} of scenarios), because its fitting time was two orders of magnitude greater than the other classes; the subset was selected by replicate order across the design rather than by outcome, so it spans every region, grid size and spatial range. Its diagnostics flagged \sdmtmbNonConvPct{}\% of fits. That figure is a property of the design: the response contains no latent spatial field for the model to estimate, so its variance components sit at zero and the package reports the boundary, as it should. The point estimates nonetheless matched the lag aware GLM's, which is the substantive result: a spatial random field is not a substitute for a lag structure, and the two address different problems. Evaluating sdmTMB where it is intended to excel, with a spatially varying effect or residual spatial dependence in the response, requires a different generating process and is the most useful extension of the repository.

\subsection{Implications for attribution and management}

Retrospective attribution remains the strongest inference that current MHW impact models support, and the reviewed literature contains defensible examples: abalone mortality in Western Australia \citep{caputi2016,hart2018}, squid range expansion in the California Current \citep{chasco2022}, pollock recruitment failure in the Gulf of Alaska \citep{rogers2021}. The simulation does not undermine those studies; it identifies the conditions under which similar studies would understate or miss an effect, and it shows that the fix is inexpensive. Forward prediction is weaker for reasons that are structural rather than methodological: responses are non stationary, lagged and mediated by trophic context, and future events will exceed the thermal range of the training data \citep{frolicher2018,welch2023}. Within those limits, the consistent biological patterns across the reviewed studies, including heightened sensitivity of early life stages, disproportionate impact of short intense events and recovery times that scale with latitude and trophic position \citep{abookire2022,guibourd2024,muhling2025}, are sufficient to support triggered harvest adjustments, temporal protections and exposure conditioned rebuilding timelines even where quantitative forecasts are not.

The geographic concentration of the evidence base is the gap that neither better methods nor better reporting will close. The methods examined here are implementable anywhere that daily temperature and an annual biological index exist, and the temperature side of that requirement is met globally by satellite products. What is missing in the tropical Pacific, the Indian Ocean and the South Atlantic is the biological series, and the regional capacity to analyse it. South to South collaboration, equitable data sharing and hosted computing environments that remove the installation barrier would do more for MHW impact science in those regions than any methodological advance.

\section{Conclusions}\label{sec:conclusions}

Marine heatwave impact assessment has matured into a field with a recognisable methodological architecture: event based exposure definitions, five families of biological response, and a hierarchy of modelling frameworks from association through structured dependence to ecosystem simulation. That architecture is uneven in practice. Association models fitted to contemporaneous, linear exposure dominate because the data and the training permit them, and this review's simulation shows that under lagged and saturating responses they understate effects, miss them entirely, and return confidence intervals that are precise about the wrong value. Distributed lag models correct both problems at modest cost; spatial random field models address a different problem and should be chosen for that reason. No single model served inference and prediction equally. The workflow, reporting standards and pitfalls table offered here are implemented in open, tested R code so that their adoption is a matter of running them rather than reconstructing them. The remaining limits to the field, sub annual biology, grid aligned response data and regional capacity, are limits of data rather than of method, and they are where investment would now return the most.
